% Options for packages loaded elsewhere
\PassOptionsToPackage{unicode}{hyperref}
\PassOptionsToPackage{hyphens}{url}
\PassOptionsToPackage{dvipsnames,svgnames,x11names}{xcolor}
\documentclass[
  10pt,
]{ctexart}
\usepackage{xcolor}
\usepackage[margin=2cm]{geometry}
\usepackage{amsmath,amssymb}
\setcounter{secnumdepth}{-\maxdimen} % remove section numbering
\usepackage{iftex}
\ifPDFTeX
  \usepackage[T1]{fontenc}
  \usepackage[utf8]{inputenc}
  \usepackage{textcomp} % provide euro and other symbols
\else % if luatex or xetex
  \usepackage{unicode-math} % this also loads fontspec
  \defaultfontfeatures{Scale=MatchLowercase}
  \defaultfontfeatures[\rmfamily]{Ligatures=TeX,Scale=1}
\fi
\usepackage{lmodern}
\ifPDFTeX\else
  % xetex/luatex font selection
  \setmainfont[]{DejaVu Sans}
\fi
% Use upquote if available, for straight quotes in verbatim environments
\IfFileExists{upquote.sty}{\usepackage{upquote}}{}
\IfFileExists{microtype.sty}{% use microtype if available
  \usepackage[]{microtype}
  \UseMicrotypeSet[protrusion]{basicmath} % disable protrusion for tt fonts
}{}
\makeatletter
\@ifundefined{KOMAClassName}{% if non-KOMA class
  \IfFileExists{parskip.sty}{%
    \usepackage{parskip}
  }{% else
    \setlength{\parindent}{0pt}
    \setlength{\parskip}{6pt plus 2pt minus 1pt}}
}{% if KOMA class
  \KOMAoptions{parskip=half}}
\makeatother
\usepackage{longtable,booktabs,array}
\usepackage{caption}
\captionsetup[table]{skip=6pt}
\newcounter{none} % for unnumbered tables
\usepackage{calc} % for calculating minipage widths
% Correct order of tables after \paragraph or \subparagraph
\usepackage{etoolbox}
\makeatletter
\patchcmd\longtable{\par}{\if@noskipsec\mbox{}\fi\par}{}{}
\makeatother
% Allow footnotes in longtable head/foot
\IfFileExists{footnotehyper.sty}{\usepackage{footnotehyper}}{\usepackage{footnote}}
\makesavenoteenv{longtable}
\setlength{\emergencystretch}{3em} % prevent overfull lines
\providecommand{\tightlist}{%
  \setlength{\itemsep}{0pt}\setlength{\parskip}{0pt}}
\usepackage{bookmark}
\IfFileExists{xurl.sty}{\usepackage{xurl}}{} % add URL line breaks if available
\urlstyle{same}
% fallback for those not using the hyperref driver hyperxmp:
\makeatletter
\@ifundefined{xmpquote}{\newcommand{\xmpquote}[1]{#1}}{}
\makeatother
\hypersetup{
  colorlinks=true,
  linkcolor={Maroon},
  filecolor={Maroon},
  citecolor={Blue},
  urlcolor={Blue},
  pdfcreator={LaTeX via pandoc}}

\author{}
\date{}

\usepackage{pdflscape,fvextra}
\setlength{\tabcolsep}{3pt}
\begin{document}

\section{LUAD transcriptomic screening
report}\label{luad-transcriptomic-screening-report}

This is a reproducible \emph{in vitro} expression screen, not a
treatment recommendation. The ranked output is in
\texttt{artifacts/\allowbreak{}reports/\allowbreak{}luad\_\allowbreak{}eh3226/\allowbreak{}all\_\allowbreak{}candidates.\allowbreak{}csv} after
running \texttt{scripts/\allowbreak{}rank\_\allowbreak{}luad\_\allowbreak{}eh3226.\allowbreak{}py}; input and output hashes
are in
\href{../data/manifests/eh3226-luad.json}{\texttt{data/\allowbreak{}manifests/\allowbreak{}eh3226-luad.\allowbreak{}json}}.
The 2.46 GB EH3226 input is fetched separately with
\texttt{scripts/\allowbreak{}fetch\_\allowbreak{}eh3226.\allowbreak{}py}.

\subsection{Prespecified method and
data}\label{prespecified-method-and-data}

GSE32863 contributes 57 verifiable LUAD tumor/normal pairs. The paired
exploratory t-test creates the disease log2 fold-change vector.
ExperimentHub EH3226 contributes A549 drug perturbation signatures at 10
µM for 24 hours. On 961 aligned landmark genes, the ranking combines
negative Spearman correlation and an up/down gene-set connectivity score
using reciprocal-rank fusion (\texttt{k=60}). In the aligned landmark
set, 57 genes meet the predetermined up threshold and 50 meet the down
threshold. No positive-control label or drug-target annotation enters
the score.

\subsection{Top-10 hypotheses}\label{top-10-hypotheses}

{\def\LTcaptype{none} % do not increment counter
\begin{longtable}[]{@{}>{\raggedright\arraybackslash}p{0.3000\textwidth}>{\raggedright\arraybackslash}p{0.3000\textwidth}>{\raggedright\arraybackslash}p{0.3000\textwidth}@{}}
\toprule\noalign{}
Rank & Drug name in EH3226 & Broad sample identity audit \\
\midrule\noalign{}
\endhead
\bottomrule\noalign{}
\endlastfoot
1 & hydrocortisone & Same connectivity; different stereochemistry \\
2 & beclomethasone-dipropionate & Exact InChIKey \\
3 & clocortolone-pivalate & Exact InChIKey \\
4 & mometasone & Identity mismatch \\
5 & flumetasone & Ambiguous GEO name \\
6 & BRD-K84203638 & No Hub sample \\
7 & hydrocortisone-hemisuccinate & Same connectivity; different
stereochemistry \\
8 & fluorometholone & No Hub sample \\
9 & fluticasone & Ambiguous GEO name \\
10 & diflorasone & No Hub sample \\
\end{longtable}
}

The \href{https://repo-hub.broadinstitute.org/repurposing}{Broad Drug
Repurposing Hub} annotates several of these names as steroids or
glucocorticoid-receptor agonists. The strong class concentration may
reflect a common A549 transcriptional response; it does not establish
therapeutic benefit in LUAD. Replaying the upstream first-duplicate
selection against official GSE92742 metadata recovers one source
\texttt{sig\_\allowbreak{}id} and \texttt{pert\_\allowbreak{}id} for every Top-10 column; see
\texttt{configs/\allowbreak{}luad\_\allowbreak{}top10\_\allowbreak{}signature\_\allowbreak{}ids.\allowbreak{}csv}. The Broad comparison
remains a separate cross-source audit, so an exact source ID does not
resolve stereochemistry, annotation mismatches, or absent Hub samples.

\subsection{Frozen-reference check}\label{frozen-reference-check}

Eleven reference names were fixed before ranking. Five exist in the
4,920-name EH3226 A549 subset: docetaxel ranks 106, crizotinib 241,
gefitinib 2,444, paclitaxel 2,962, and erlotinib 3,147. The other six
are absent and are \textbf{unmeasured}, not failures. The spread of
these ranks and the top-ranked steroid cluster mean this screen has not
demonstrated LUAD drug recovery. No AUC was calculated by treating
unmeasured or unknown drugs as negatives.

\subsection{Evidence and limitations}\label{evidence-and-limitations}

All ten candidates retain \texttt{insufficient\_\allowbreak{}evidence} in their
structured ledgers. A frozen
\href{https://www.ncbi.nlm.nih.gov/books/NBK25499/}{NCBI PubMed
E-utilities} title/abstract search for each exact drug name with
LUAD/NSCLC terms returned 39 hits for hydrocortisone, one for
beclomethasone dipropionate, and zero for the other eight names. This
narrow query is a retrieval snapshot, not an evidence-absence test;
synonyms and broader terms were not searched. Two records were checked
at the abstract level and placed in \texttt{context\_\allowbreak{}evidence} only:
\href{https://pubmed.ncbi.nlm.nih.gov/35676421/}{PMID 35676421}
independently proposes hydrocortisone as a computational NSCLC
candidate, without efficacy validation in its abstract;
\href{https://pubmed.ncbi.nlm.nih.gov/12538830/}{PMID 12538830} reports
glucocorticoid-receptor-dependent CYP3A5 induction by beclomethasone
dipropionate in A549 cells, a mechanistic observation rather than
anticancer efficacy. Other retrieved hits remain unreviewed.

The expanded \href{luad_top10_evidence_matrix.md}{Top-10 identity and
evidence matrix} resolves the dominant glucocorticoid class, source
signature IDs, name-level targets and pathways, candidate-specific
context for hydrocortisone and beclomethasone dipropionate, and shared
supporting/contradicting evidence. It does not upgrade any candidate: no
compound has independently validated LUAD efficacy, dose relevance,
safety, or patient-context evidence. The
\href{luad_experimental_validation_protocol.md}{experimental validation
protocol} prespecifies dose, molecular-subtype, mechanism, and
combination tests, but no wet-lab measurements are claimed. The current
expression rank is useful for prioritizing research only.

\end{document}
